\hypertarget{under-the-hood-built-in-data-structures}{%
\chapter{UNDER THE HOOD: BUILT-IN DATA
STRUCTURES}\label{under-the-hood-built-in-data-structures}}

\hypertarget{python-list-memory-management}{%
\subsection{5.1 Python List Memory
Management}\label{python-list-memory-management}}

A Python list is a contiguous dynamic array of
\passthrough{\lstinline!PyObject*!} pointers. Because list sizes change
dynamically, CPython overallocates memory blocks to achieve \(O(1)\)
amortized append performance.

\hypertarget{the-list-growth-formula}{%
\subsubsection{The List Growth Formula:}\label{the-list-growth-formula}}

When a list exceeds its current capacity during an
\passthrough{\lstinline!append!} or \passthrough{\lstinline!insert!}
call, CPython recalculates its allocation size in
\passthrough{\lstinline!Objects/listobject.c!} using the formula:
\[\text{new\_allocated} = \text{newsize} + (\text{newsize} \gg 3) + (\text{newsize} < 9 \,?\, 3 : 6)\]

Here is the resulting capacity growth trace: \textbar{} Item Count
(\passthrough{\lstinline!newsize!}) \textbar{} Allocated Capacity
\textbar{} Overallocation Factor \textbar{}
\textbar---\textbar---\textbar---\textbar{} \textbar{} 0 \textbar{} 0
\textbar{} - \textbar{} \textbar{} 1 \textbar{} 4 \textbar{} \(400\%\)
\textbar{} \textbar{} 5 \textbar{} 8 \textbar{} \(160\%\) \textbar{}
\textbar{} 9 \textbar{} 16 \textbar{} \(177\%\) \textbar{} \textbar{} 17
\textbar{} 25 \textbar{} \(147\%\) \textbar{} \textbar{} 1000 \textbar{}
1129 \textbar{} \(112.9\%\) \textbar{}

This formula balances memory overhead and reallocation speed. By adding
\(\approx 12.5\%\) extra slots as the list grows, it minimizes heap
reallocations and memory copying costs.

\hypertarget{python-dict-hashing-and-collision-pertubation}{%
\subsection{5.2 Python Dict: Hashing and Collision
pertubation}\label{python-dict-hashing-and-collision-pertubation}}

A Python dictionary is a sparse hash table. CPython uses open addressing
(probing) to resolve collisions.

\hypertarget{collision-perturbation-formula}{%
\subsubsection{Collision Perturbation
Formula:}\label{collision-perturbation-formula}}

When a hash collision occurs, CPython searches for the next index using
a pseudo-random probing formula. Rather than a simple linear probe
(which causes clustering), it shifts the upper bits of the hash code
down to contribute to the index calculation:
\[i = (5 \times i + 1 + \text{perturb}) \pmod{\text{size}}\]

In each probe step: \[\text{perturb} \gg= 5\] As \(\text{perturb}\)
drops to zero, the formula collapses to:
\[i = (5 \times i + 1) \pmod{\text{size}}\] This ensures every slot in
the table is eventually checked.

\hypertarget{the-compact-dict-layout-pep-468-python-3.6}{%
\subsubsection{The Compact Dict Layout (PEP 468 / Python
3.6+):}\label{the-compact-dict-layout-pep-468-python-3.6}}

Previously, dictionary tables were stored as sparse arrays of
\passthrough{\lstinline!PyDictEntry!} structs (24 bytes per bucket):

\begin{lstlisting}
Pre-3.6 Sparse Dictionary Layout (Huge footprint, unpreserved order):
[ Index 0: <hash, key_ptr, value_ptr> (24 bytes) ]
[ Index 1: <Empty / 24 bytes null padding>       ]
[ Index 2: <hash, key_ptr, value_ptr> (24 bytes) ]
\end{lstlisting}

The modern compact dict layout splits this into two arrays: 1.
\textbf{Indices Array}: A sparse array of small integers (typically 1
byte each) mapping hash codes to positions in the entries array. 2.
\textbf{Entries Array}: A dense, contiguous array of
\passthrough{\lstinline!PyDictKeyEntry!} structs (containing
\passthrough{\lstinline!hash!}, \passthrough{\lstinline!key\_ptr!}, and
\passthrough{\lstinline!value\_ptr!}) stored in insertion order.

\begin{lstlisting}
Modern Compact Dictionary Layout (PEP 468):
Indices Array: [ 0, -1, 1 ]  (Fast index references)
Entries Array (Dense, ordered):
[ Slot 0: <hash1, key1, val1> ]
[ Slot 1: <hash2, key2, val2> ]
\end{lstlisting}

This saves up to 40\% memory and preserves insertion order by default.

\hypertarget{sets-and-tuples-optimizations}{%
\subsection{5.3 Sets and Tuples
Optimizations}\label{sets-and-tuples-optimizations}}

\begin{itemize}
\tightlist
\item
  \textbf{Sets (\passthrough{\lstinline!PySetObject!})}: Implemented
  similarly to compact dicts, but contain only keys. Set lookups bypass
  value lookups entirely, using specialized open-addressing C loops.
\item
  \textbf{Tuples (\passthrough{\lstinline!PyTupleObject!})}: Tuples are
  immutable, so they are allocated as a single contiguous block of
  memory containing both the \passthrough{\lstinline!PyVarObject!}
  header and the array of \passthrough{\lstinline!PyObject*!} pointers.

  \begin{itemize}
  \tightlist
  \item
    \emph{Allocation Cache}: CPython avoids calling the system memory
    allocator when creating short tuples. It maintains a \textbf{free
    list} of recycled tuple objects for sizes up to 20. When a tuple is
    garbage-collected, its memory is not released to the OS; instead, it
    is placed in the free list for quick reuse.
  \end{itemize}
\end{itemize}

\begin{center}\rule{0.5\linewidth}{0.5pt}\end{center}
